  \subsubsection{余りと周期}\label{節：ガウス型の余りと周期}
    整数$X$と正整数$m$に対して,
    \[
      X-m\left\lfloor\frac Xm\right\rfloor
    \]
    は$X$を$m$で割った余りであり,\,$0,1,\ldots,m-1$のいずれかである.
    \bookterm{floor-function}{床関数}の定義から,\textbf{値は0以上$m$未満の整数に限られる}.

    \,$x\equiv y\pmod m$は,\,$x-y$が$m$の倍数であることを表す合同式である.
    次の項が$n$にも依存する場合は,$a_n$の値だけでなく,$n$を$m$で割った余りも一緒に記録する.
    同じ組が再び現れ,その組から次の組が一意に決まるなら,そこから先は同じ並びを繰り返す.ただし,初項から周期的とは限らない.

    \noindent \bookblocklink{example-gauss-role-3}{problem}{answer}{【例題】}\par\nobreak
    \begin{enumerate}[label=(\arabic*),parsep=3pt,beginpenalty=10000,format=\bookitemlink{example-gauss-role-3}{problem}{answer}]
      \item \bookterm{sequence}{数列}を次のように定める.\bookterm{general}{一般項}と最小の正の周期を求めよ.
            \[
              a_1=1,\qquad
              a_{n+1}=X_n-4\left\lfloor\frac{X_n}{4}\right\rfloor,
            \]
            \[
              X_n=a_n^2+n^2+n\quad(n\ge1).
            \]
    \end{enumerate}

    \noindent \bookblocklink{example-gauss-role-3}{hint}{problem}{【方針】}\par\nobreak
    \begin{enumerate}[label=(\arabic*),parsep=3pt,beginpenalty=10000,format=\bookitemlink{example-gauss-role-3}{hint}{problem}]
      \item 次の項は$X_n$を4で割った余りです.
            \,$(a_n,n\bmod4)$の組を追うと,同じ状態への戻りを確かめられます.
    \end{enumerate}

    \noindent \bookblocklink{example-gauss-role-3}{answer}{problem}{【解答】}\par\nobreak
    \begin{enumerate}[label=(\arabic*),parsep=3pt,beginpenalty=10000,format=\bookitemlink{example-gauss-role-3}{answer}{problem}]
      \item \[
              a_n=
              \begin{cases}
                1 & (n\equiv0,1\pmod4),\\
                3 & (n\equiv2,3\pmod4).
              \end{cases}
            \]
            最小の正の周期は4である.
    \end{enumerate}

    \noindent \bookblocklink{example-gauss-role-3}{solution}{problem}{【解法】}\par\nobreak
    \begin{enumerate}[label=(\arabic*),parsep=3pt,beginpenalty=10000,format=\bookitemlink{example-gauss-role-3}{solution}{problem}]
      \item 初項と更新から各項は整数で,第2項以降は$0\le a_n<4$である.
            \[
              a_{n+1}\equiv a_n^2+n(n+1)\pmod4
            \]
            となる.順に追うと,
            \[
              a_1,a_2,a_3,a_4,a_5=1,3,3,1,1.
            \]
            第1項と第5項では$(a_n,n\bmod4)=(1,1)$が一致する.
            同じ組から同じ次の組が決まるので,帰納的に$a_{n+4}=a_n$となり,【解答】を得る.
            第4項は値だけが同じで,添字の余りが違う.

            $a_1\ne a_2,a_3$と$a_2\ne a_5$より,1,2,3は周期でない.従って最小周期は4である.
    \end{enumerate}
